% Document Setup ---------------------------------------------------------

% Set document class.
\documentclass[10pt]{article}

% Set document details.
\title{\Large\textsc{Frontiers in Computing: Discussion 1}}
\author{\large\textsc{Rento Saijo}}
\date{2026-09-06}

% Load packages.
\usepackage[letterpaper, margin = 1in]{geometry}
\usepackage{amsmath, amssymb}
\usepackage{enumitem}
\usepackage{graphicx}
\usepackage{fancyhdr, lastpage}
\usepackage{tcolorbox}
\usepackage[hidelinks]{hyperref}

% Page Layout ------------------------------------------------------------

% Configure equations and paragraphs.
\allowdisplaybreaks
\setlength{\parindent}{0pt}
\setlength{\parskip}{0.3em}
\clubpenalty = 10000
\widowpenalty = 10000

% Configure page footer.
\pagestyle{fancy}
\fancyhf{}
\fancyfoot[C]{\textsc{Page \thepage\ of \pageref{LastPage}}}
\renewcommand{\headrulewidth}{0pt}

% Helpers ----------------------------------------------------------------

% Define problem environment.
\newtcolorbox{problem}[1][]{
    colframe = darkgray,
    arc      = 0.1mm,
    boxrule  = 2pt,
    boxsep   = 1mm,
    title    = {\textsc{#1}}
}

% Content ----------------------------------------------------------------

% Render document title.
\begin{document}
\maketitle
\thispagestyle{fancy}
\vspace{-1em}\rule{\linewidth}{0.6pt}\vspace{0.5em}

% Write problem statement.
\begin{problem}[Problem 1]
    Summarize the videos.
\end{problem}

% Introduce integrated discussion.
An email service, a video platform, and a county social-services department all depend on computing resources that must keep working as their users' needs change. Behind each service are questions about where information lives, who maintains the equipment, how much capacity is available, and who can access it. Taken together, the videos explain cloud computing through these practical concerns. Providers make computing capacity available as services, while customers choose how much of the associated work to entrust to them. Understanding the cloud therefore requires us to connect its physical infrastructure with the decisions organizations make about responsibility, cost, performance, and security.

% Explain physical infrastructure and resource sharing.
\section*{From local computers to cloud infrastructure}

We can begin with a familiar computer. Its processor and memory provide compute capacity, its drives store applications and data, and its network connections allow it to communicate. An operating system makes those capabilities usable, and applications turn them into activities we care about, such as writing documents or running a business. An organization's data center extends this arrangement across many servers, storage systems, and network devices. The public cloud draws on the same physical foundations: providers operate large collections of equipment in data centers and expose their capabilities through services that customers access over networks.

The email example makes this arrangement concrete. An organization that runs its own email server must obtain hardware, install an operating system and email software, configure the system, and maintain it. A failed component or software problem becomes its responsibility to resolve. With a hosted email service, much of that operational work moves to the service provider. Similarly, a video creator who uses YouTube gains access to an existing platform for storing and delivering videos. Building an equivalent service independently would require a server, supporting software, and substantial network bandwidth. YouTube illustrates how outsourcing computing work can be useful even when the commercial arrangement involves advertising revenue instead of a direct server-rental bill.

An important step between physical equipment and cloud services is virtualization. A \textit{hypervisor} creates virtual machines, each with an allocation of processor capacity, memory, storage, and networking. Multiple virtual machines can run on one physical server, and each can have its own operating system and application environment. One might run Windows, while another runs Linux with a different set of libraries. The hypervisor provides isolation and controls resource allocation, allowing these workloads to share hardware while maintaining separate execution environments. Organizations can consequently use equipment more efficiently and create additional machines without waiting for new physical servers to be purchased, delivered, and installed.

Containers introduce another form of isolation. As Savill explains, virtual machines virtualize hardware, while containers provide isolated environments within an operating system. Sharing that operating system reduces the overhead associated with running a separate operating system for every workload and allows containers to start quickly. Kubernetes helps manage these environments, including their creation and health. These technologies support the broader idea of resource pooling: a provider combines physical capacity and allocates portions of it to different customers and workloads. In a public cloud, this commonly produces a \textit{multitenant} environment, where customers share underlying infrastructure and rely on isolation mechanisms to keep their resources separate.

% Compare service responsibilities and deployment choices.
\section*{Choosing service and deployment models}

Cloud service models describe how responsibilities are divided between a customer and a provider. In an entirely on-premises arrangement, the organization manages the whole stack: networking, storage, servers, virtualization, operating systems, supporting software, applications, and data. Moving to a cloud service changes that division of work. The central choice is how much control the organization needs and how much maintenance it is prepared to perform.

With \textit{infrastructure as a service} (IaaS), the provider manages the physical infrastructure and virtualization layer. A virtual machine in the cloud is the clearest example. The customer chooses and maintains its operating system, installs the necessary software, and manages the applications and data inside it. This flexibility is valuable when an application requires a particular operating system, specific libraries, or direct access to system settings. It also carries responsibilities for patching, configuration, protective software, and planning backup and recovery for the workload. Supporting cloud tools may help perform these tasks, but the customer still has to ensure that they are addressed.

With \textit{platform as a service} (PaaS), the provider also manages the operating system and relevant runtime environment. The customer concentrates on its own application and associated data. For a business developing a custom web application, this can remove much of the work involved in maintaining the underlying platform. Savill also discusses serverless computing, in which an event triggers a piece of code or a workflow. The trigger could be a schedule, an incoming request, or an event in another service. Servers still perform the computation; the service takes on more of their management so that the developer can focus on the code and its purpose.

With \textit{software as a service} (SaaS), the provider supplies the application itself. Google Docs, for example, lets users create and edit documents through a browser without installing and maintaining their own office-application service. Microsoft 365 and customer-management applications follow the same general model. The provider maintains the software and supporting infrastructure, while the organization administers how people use the service, including their access and roles. This arrangement substantially reduces maintenance work, although customers still have responsibilities concerning their identities, devices, and use of information.

Savill's comparison between a general-purpose PC and a game console helps explain the tradeoff. A PC supports many kinds of software and gives its owner considerable flexibility, alongside maintenance duties involving the operating system, drivers, and security. A console provides a more specialized experience in which the manufacturer manages much of that complexity. The appropriate choice depends on what we want to accomplish. A ready-made SaaS product can serve an existing business need; PaaS supports an organization building its own application; and IaaS remains useful when detailed system control is necessary. The service models offer different ways to match responsibilities to requirements.

Deployment models answer a separate question: who uses the cloud infrastructure and how it is organized. A \textit{public cloud} serves many customers through a provider such as AWS, Microsoft Azure, or Google Cloud. A \textit{private cloud} is dedicated to one organization and provides cloud capabilities for its own use. A \textit{community cloud} serves a group with shared requirements, such as organizations working within a common institutional setting. A \textit{hybrid cloud} combines cloud environments, allowing workloads to operate across them. Savill emphasizes that an organization may retain systems locally because of a connection to an existing mainframe, location requirements, or the absence of a suitable provider region. Service and deployment choices therefore work together: an organization must decide both where its workloads belong and which parts of their operation it wants a provider to manage.

% Explain elasticity and consumption-based costs.
\section*{Elasticity, self-service, and cloud economics}

The videos identify several characteristics that make cloud services distinctive. Resource pooling supplies shared capacity; broad network access makes services reachable through network connections; and on-demand self-service lets authorized users obtain resources when they need them. Rapid elasticity allows capacity to expand and contract, while measured service tracks consumption. Together, these characteristics change how organizations acquire and pay for computing resources.

Consider Savill's example of a web service that needs ten running instances during a busy period but only three when demand falls. Horizontal scaling adjusts the number of instances serving the workload. An organization can schedule changes around predictable patterns or use autoscaling to respond to changing demand. A pizza business might anticipate a rush on Friday evenings or during the Super Bowl, while a tax-related service experiences a different seasonal pattern. When traffic decreases, removing unnecessary instances brings the deployed capacity back toward what the application currently needs. This ability to contract is as useful economically as the ability to expand is operationally.

Measured use connects those resource decisions to expenditure. Depending on the service, consumption may be represented by virtual-machine running time, storage capacity, database activity, or the input and output tokens processed by an AI service. Consequently, organizations need to understand what is being measured for the services they select. The opportunity to pay according to use becomes most valuable when they actively adjust what they run. Leaving excessive capacity in place continues to incur costs even when the application receives little benefit from it.

This differs from purchasing equipment to cover several years of expected growth. A traditional deployment may require a substantial initial investment in servers and storage before the future workload is known. A startup using cloud services can begin with a small allocation and increase it as the business develops. If a project ends, it can release the resources it no longer needs. Savill extends the example to a workload requiring a powerful GPU-based system for only an hour each month. Temporary access makes such a workload feasible without purchasing a machine that would sit unused for most of the month.

The same principle helps explain the appeal of outsourcing at a larger scale. PowerCert uses Netflix's use of AWS for workloads such as databases, analytics, and video transcoding to illustrate how a company can obtain extensive computing services and devote more attention to its business. However, flexible access still requires technical planning and oversight. Organizations must choose appropriate services, manage their consumption, and establish controls over who can create resources. Cloud economics follows from those choices as well as from the provider's billing model.

% Connect architecture, storage, and managed services.
\section*{Performance, availability, and storage}

Once resources can be created on demand, their placement becomes an important design decision. Providers organize infrastructure into geographic \textit{regions}. Within a region, availability zones provide groups of facilities with independent supporting systems, such as power, cooling, and networking. Distributing an application across zones can help it continue operating when one location experiences a failure. For disruption affecting a whole region, an organization may also design a deployment that uses another region. Savill distinguishes an active-active arrangement, where services already run in multiple locations, from an arrangement in which a secondary deployment is available to start when needed. These options make resilience possible, but the application must be configured to use them.

Geographic distribution also addresses latency, the time required for information to travel between locations. A fraction of a second can appear negligible to a person, yet become significant when an application and its database exchange many messages to complete one operation. Repeated delays accumulate into an experience that users perceive as slow or unresponsive. Locating services close to their users and considering the distance between interacting components can reduce this problem. Thus, regions serve two related purposes: they provide choices for recovering from geographically concentrated failures and help place services near the people and systems that use them.

Storage introduces further choices. Block storage presents data in addressable blocks, while file storage organizes information into shared files and folders. Object storage is suited to large collections of unstructured content, including images and videos. These forms support different access patterns, so selecting storage requires attention to how an application reads and changes its data. Savill's examples of network file services and blob storage illustrate that the cloud offers several ways to expose stored information, beyond simply attaching a disk to a virtual machine.

Microsoft also distinguishes hot storage for frequently accessed information from colder options for information used less often or retrieved less urgently. Storage performance depends on the underlying media and the network connecting the user to the service. Choosing faster media such as SSDs can improve performance, with a corresponding cost consideration. These decisions connect the application's needs to the characteristics of the service supporting it.

Content delivery networks (CDNs) apply geographic distribution to content delivery. Microsoft's tourism example considers videos promoting the United Arab Emirates: distributing the content closer to viewers can improve their experience across locations. This extends the latency principle to the information an audience requests. Distribution also supports availability and the ability to serve greater demand.

Beyond storage and virtual machines, Savill describes a broad collection of managed capabilities. These include databases, caching, analytics, search, AI and machine learning, and services through which connected devices exchange information. Vector search provides an example of how these services can support new application behavior. A model represents text, images, or audio numerically, and a search can use the distance between representations to identify similar content. This supports matching that does not require identical wording. The cloud's scale also makes substantial computing resources available for demanding AI workloads. Organizations can combine these capabilities with their own applications, expanding what they can build without independently operating every underlying system.

% Explain shared security responsibilities and governance.
\section*{Security and governance}

Security runs through every earlier decision because the division of operational work also determines who must protect each part of a system. Providers invest in securing their facilities and services. Customers, in turn, must address the responsibilities associated with the resources they choose. An organization running IaaS virtual machines retains work inside those machines; one using SaaS relies on the provider for application maintenance while still controlling how its users access and handle information. Backup and recovery similarly need to be understood in relation to the selected service and the customer's workload. A resilient provider infrastructure forms part of a reliable solution, alongside the way the organization configures and operates its resources.

Microsoft presents \textit{defense in depth} as a series of protective layers organized around confidentiality, integrity, and availability. Physical safeguards restrict access to facilities, while identity controls govern access to resources. Perimeter defenses address distributed denial-of-service attacks that attempt to overwhelm a service. Network controls limit communication, and protections within computing systems and applications address vulnerabilities. Data access controls complete this layered approach, while encryption helps protect stored information and communications. The layers work together so that protecting information involves the path to it as well as the system holding it.

Identity receives particular attention in Savill's explanation. A cloud identity provider can support a single identity across multiple services, giving users a single sign-on experience. An organization with an existing directory can also synchronize identities to a cloud provider. This simplifies access, while concentrating the need to protect the identity that opens those services. Multifactor authentication, passwordless methods, and policies informed by risk signals help secure that access. For example, attempted logins from locations too far apart for a person to travel between in the available time may prompt stronger authentication or additional restrictions. Device health can also inform an access decision.

These practices reflect a \textit{zero-trust} approach, which continually evaluates access instead of assuming that everything inside a network is trustworthy. An attacker who reaches a network should still encounter controls protecting individual resources and information. The organization therefore needs to know who is requesting access, what they are trying to do, and whether the surrounding conditions are acceptable. Centralized identity management makes it possible to monitor these decisions across services and respond to signs of risk.

Broad network access requires a further distinction between the \textit{control plane} and the \textit{data plane}. The control plane provides management operations, such as creating, changing, or deleting a resource through a portal, command-line interface, or API. Identity and role-based permissions protect those operations. The data plane concerns access to what the resource actually provides; for a virtual machine, Savill uses an SSH or remote-desktop connection as an example. Choosing a public cloud does not require exposing every such connection directly to the Internet. Private connections, encrypted tunnels between locations, firewalls, and time-limited access can constrain how resources are reached.

Governance brings these controls into everyday resource management. In a traditional workflow, an IT team might review each request before creating a resource. With controlled self-service, the team can establish policies in advance and let authorized users work within them. Policies can restrict resource types and locations, prevent inappropriate public exposure, and limit which actions particular roles can perform. Quotas and oversight also help manage consumption. This preserves organizational requirements while giving teams quicker access to useful services. Security and flexibility can therefore be designed into the same operating process.

% Connect technical choices to organizational outcomes.
\section*{Applying cloud capabilities to organizational needs}

The Microsoft series connects these concepts through the experience of Chesterfield County's Department of Social Services in Virginia. The department had relied on longstanding infrastructure and fragmented information, with staff using spreadsheets and handwritten records to support their work. Collecting information across sources was difficult, and manual processes contributed to duplicate entries and mistakes. Increased demand for assistance added pressure to an already burdensome workflow. The case gives the preceding technical discussion a concrete purpose: employees needed a more effective way to organize information and deliver services.

The department adopted Microsoft Dynamics 365, using a SaaS application to centralize data and document management and streamline its processes. Microsoft reports a 33\% reduction in the time needed to complete common tasks, together with a reduced likelihood of documents being duplicated or lost. The series also describes changes to employee training and the department's continued use of related digital tools. These outcomes illustrate how an application and its supporting processes can translate cloud capabilities into improvements in daily work.

The case also brings the choice of service model into focus. The department's immediate need concerned information management and service delivery, making access to a suitable application more useful than simply acquiring additional server capacity. Other organizations may need a platform for custom development or virtual machines for existing software. In each case, the design must connect the selected services to the work being done. Requirements for growth, day-to-day operation, migration, security, and integration with existing systems all help shape that design, and the arrangement may evolve as the organization changes.

Across the videos, the cloud emerges as a way to make computing capacity and specialized capabilities broadly accessible. Its value comes from connecting that access to an appropriate division of responsibility. Virtualization and pooling support flexible resources; service models determine how much of their operation the provider manages; geographic design and storage choices shape performance; and identity and governance guide their use. When these decisions reinforce the organization's purpose, cloud services can free time and resources for the activities that matter to its users, from sharing a document to delivering public assistance.

% Identify source videos.
{\small
\textbf{Source videos.} PowerCert Animated Videos, \textit{Cloud Computing Explained}; John Savill's Technical Training, \textit{Introduction to Cloud}; and Microsoft, \href{https://wwps.microsoft.com/blog/intro-cloud-computing-training-video/}{\textit{Introduction to Cloud Computing}}, including the modules on cloud computing, service and deployment models, security, storage, and cloud design.
\par}

% End document.
\end{document}
